\begin{center}\begin{tikzpicture}[font=\small,>=stealth,line join=round]
\draw[gray!40,rounded corners=3pt] (0,0) rectangle (14.2,8.3);\node at (7.1,7.700000000000001) {负指数：越靠近轴，不等于真的碰到轴};
\begin{scope}[shift={(1.6,2.033333333333333)},x=1.0cm,y=0.7666666666666666cm]\draw[->,gray] (-0.5,0)--(5.5,0) node[right]{$x$};\draw[->,gray] (0,-0.5)--(0,5.5) node[above]{$y$};\draw[orange!85!black,thick,solid,domain=0.2:5,samples=180] plot(\x,{1/\x});\fill[orange!85!black] (0.25,4) circle(2.2pt);\fill[orange!85!black] (4,0.25) circle(2.2pt);\node[right=5pt] at (.25,4){$(\frac14,4)$};\node[above=5pt] at (4,.25){$(4,\frac14)$};\end{scope}\node[anchor=west,align=left,text width=5.0cm] at (8.4,4.0) {$y=1/x,\ x>0$\\[9pt]$(\frac14,4)$与$(4,\frac14)$\\[9pt]$x\to0^+$：$y\to+\infty$\\[9pt]$x\to+\infty$：$y\to0^+$\\[9pt]始终有$x>0,y>0$};
\node[text width=12.6cm,align=center] at (7.1,.65) {图取负指数中的1/x；一般x的负p次方具有相同的正半轴渐近行为，p为正数。};
\end{tikzpicture}\end{center}
